%!TEX root = ../dokumentation.tex

\chapter{Projektumfang (Project Scope)}

\section{Scope Statement}

Ziel des Projekts ist die Einführung einer konzernweit einheitlichen Lernplattform für die dm-Gruppe. Die Plattform unterstützt Planung, Buchung, Durchführung, Dokumentation sowie Auswertung von Schulungs- und Weiterbildungsmaßnahmen und ersetzt bestehende dezentrale Lern- und Schulungssysteme.

Die Lernplattform bildet eine integrierte End-to-End-Lösung für das konzernweite Weiterbildungsmanagement. Lerninhalte, Nutzerverwaltung sowie Weiterbildungs- und Kompetenzdaten werden zentral konsolidiert und bilden eine organisationsweite \textit{Single Source of Truth} für Reporting und Steuerung.

Das Projekt umfasst sowohl die technische Einführung der Plattform als auch die Standardisierung der zugrunde liegenden Prozesse sowie die organisatorische Einführung im Unternehmen. Der Projektumfang orientiert sich am definierten \ac{MOV}, insbesondere der Reduktion administrativer Aufwände im Schulungsmanagement sowie der Bereitstellung organisationsweiter Weiterbildungs- und Kompetenzkennzahlen.

\section{In Scope}

Die folgenden Leistungen und Deliverables sind Bestandteil des Projekts:

\begin{itemize}
    \item Auswahl und Einführung einer zentralen Lernplattform (z.\,B. SAP SuccessFactors Learning)
    \item Konfiguration der Plattform gemäß organisatorischen Anforderungen
    \item Einrichtung der Systemlandschaft (Produktiv-, Entwicklungs- und Testumgebungen)
    \item Analyse bestehender Schulungs- und Weiterbildungsprozesse
    \item Definition und Implementierung konzernweiter Soll-Prozesse für Planung, Buchung, Durchführung, Dokumentation und Reporting von Schulungen
    \item Migration relevanter Stamm-, Kurs- und Historien-Daten aus bestehenden Systemen
    \item Aufbau eines zentralen Reporting- und Analysemodells für Weiterbildungs- und Kompetenzkennzahlen
    \item Definition eines konzernweiten Rollen- und Berechtigungskonzepts
    \item Schulung von Administratoren und Fachbereichen
    \item Kommunikations- und Change-Management-Maßnahmen zur Einführung der Plattform
    \item Begleitung des konzernweiten Rollouts
\end{itemize}

\section{Out of Scope}

Die folgenden Themen sind nicht Bestandteil des Projekts:

\begin{itemize}
    \item Entwicklung individueller Fachbereichs-Sonderlösungen
    \item Entwicklung einer eigenen Lernplattform (Individualsoftware)
    \item Parallelbetrieb bestehender Lernsysteme nach Abschluss der Migration
    \item Erstellung oder Pflege spezifischer Lerninhalte durch das Projektteam
    \item Anpassung von \ac{HR}- oder Bonusmodellen selbst (die Plattform stellt lediglich die notwendige Datenbasis bereit)
\end{itemize}

\section{Key Deliverables}

Die wichtigsten Projektergebnisse sind:

\begin{itemize}
    \item Produktiv eingeführte zentrale Lernplattform
    \item Standardisierte konzernweite Schulungs- und Weiterbildungsprozesse
    \item Migration und Konsolidierung relevanter Weiterbildungsdaten
    \item Zentrales Reporting- und \ac{KPI}-System für Weiterbildung und Kompetenzen
    \item Definiertes Governance- und Rollenmodell für den Betrieb der Plattform
    \item Schulungs- und Rolloutkonzept für die organisatorische Einführung
\end{itemize}
